\documentclass[UTF8,a4paper,11pt]{ctexart}

\usepackage[a4paper,margin=2.4cm]{geometry}
\usepackage{booktabs}
\usepackage{tabularx}
\usepackage{longtable}
\usepackage{array}
\usepackage{xcolor}
\usepackage{enumitem}
\usepackage{hyperref}
\usepackage{amsmath,amssymb}
\usepackage{tikz}
\usepackage[most]{tcolorbox}
\usepackage{titlesec}
\usepackage{fancyhdr}
\usepackage{caption}
\usepackage{float}
\providecommand{\tightlist}{\setlength{\itemsep}{0pt}\setlength{\parskip}{0pt}}
\newcounter{none}

\definecolor{mainblue}{RGB}{28,90,160}
\definecolor{lightblue}{RGB}{235,244,255}
\definecolor{darkgray}{RGB}{70,70,70}
\definecolor{accentorange}{RGB}{230,120,40}
\definecolor{softgreen}{RGB}{235,250,240}

\hypersetup{colorlinks=true,linkcolor=mainblue,urlcolor=mainblue,citecolor=mainblue}

\titleformat{\section}{\Large\bfseries\color{mainblue}}{\thesection}{0.8em}{}
\titleformat{\subsection}{\large\bfseries\color{darkgray}}{\thesubsection}{0.8em}{}
\titleformat{\subsubsection}{\normalsize\bfseries\color{darkgray}}{\thesubsubsection}{0.8em}{}

\pagestyle{fancy}
\setlength{\headheight}{14pt}
\fancyhf{}
\lhead{选题前置调研}
\rhead{研究方案报告}
\cfoot{\thepage}

\tcbset{
  mybox/.style={colback=lightblue,colframe=mainblue,arc=2mm,boxrule=0.8pt,left=2mm,right=2mm,top=1.5mm,bottom=1.5mm},
  keybox/.style={colback=softgreen,colframe=green!50!black,arc=2mm,boxrule=0.8pt,left=2mm,right=2mm,top=1.5mm,bottom=1.5mm},
  warnbox/.style={colback=orange!8,colframe=accentorange,arc=2mm,boxrule=0.8pt,left=2mm,right=2mm,top=1.5mm,bottom=1.5mm}
}

\title{\vspace{-1.5cm}\textbf{选题前置调研报告}\\[0.5em]\large TODO：项目副标题或英文标题}
\author{TODO}
\date{\today}

\begin{document}
\maketitle

\begin{tcolorbox}[mybox,title=一句话概括]
TODO：用 2--4 句话说明项目到底研究什么、不是研究什么、为什么现在值得做、目前可行性是 Pass / Weak Pass / Fail。
\end{tcolorbox}

\tableofcontents
\newpage

\section{项目核心判断}

\subsection{核心问题不是}

TODO。

\subsection{真正值得研究的问题}

\begin{tcolorbox}[keybox,title=核心研究问题]
TODO。
\end{tcolorbox}

\subsection{本项目定位}

\begin{table}[H]
\centering
\caption{本项目定位}
\begin{tabularx}{\textwidth}{p{3cm}X}
\toprule
\textbf{维度} & \textbf{定位} \\
\midrule
不是 & TODO \\
是 & TODO \\
核心输入 & TODO \\
核心输出 & TODO \\
核心价值 & TODO \\
\bottomrule
\end{tabularx}
\end{table}

\section{为什么适合目标会议或期刊}

TODO。

\section{核心概念与研究问题}

TODO。

\section{输入模态、数据模态或研究对象边界}

TODO。

\section{系统方案或方法方案}

TODO。

\section{交互/方法设计原则}

TODO。

\section{算法与数据的最小可行方案}

TODO。

\section{典型案例或使用场景}

\begin{tcolorbox}[mybox,title=Demo Case]
\textbf{场景输入：} TODO\\
\textbf{系统/方法判断：} TODO\\
\textbf{推荐输出：} TODO\\
\textbf{评价信号：} TODO
\end{tcolorbox}

\section{研究问题与评价方案}

TODO。

\section{数据集支持与可行性判断}

\begin{tcolorbox}[warnbox,title=可行性结论]
\textbf{Verdict：} TODO：Pass / Weak Pass / Fail。
\end{tcolorbox}

\section{标题润色方案与摘要}

TODO。

\section{预期论文结构}

TODO。

\section{投稿定位}

TODO。

\section{风险与应对}

TODO。

\section{下一步工作}

TODO。

\section{最终总结}

\begin{tcolorbox}[keybox,title=最终总结]
TODO：用一段强判断收束，说明这个方向应该继续、收缩还是放弃。
\end{tcolorbox}

\end{document}

